\section{三角系统}

在这一节中，我们只讨论两类非常简单的问题
\begin{Equation}[三角系统]
    \vb{L}\vb{x}=\vb{b}\qquad \vb{U}\vb{x}=\vb{b}
\end{Equation}

其中，$\vb{L}\in\R^{n\times n}$为下三角矩阵，$\vb{U}\in\R^{n\times n}$为上三角矩阵。

\subsection{前向代换}

考虑一个$2\times 2$的下三角系统
\begin{Equation}[2x2前向代换1]
    \mqty[
        l_{11} & 0      \\
        l_{21} & l_{22} \\
    ]
    \mqty[
        x_1 \\
        x_2 \\
    ]
    =
    \mqty[
        b_1 \\
        b_2 \\
    ]
\end{Equation}

若$l_{11},l_{22}\neq 0$，则未知数可以从上到下依次确定
\begin{Gather}[2x2前向代换]
    x_1=b_1/l_11\id{2}\\
    x_2=(b_2-l_{21}x_1)/l_{22}\id{3}
\end{Gather}

这一问题的一般形式可以概括为
\begin{BoxFormula}[前向代换]
    若$\vb{L}\in\R^{n\times n}$是下三角矩阵，而$\vb{b}\in\R^n$，那么
    \begin{Equation}[下三角系统]
        \vb{L}\vb{x}=\vb{b}
    \end{Equation}

    该方程可以用前向代换（Forward Substitution）的方式求出
    \begin{Equation}[前向代换]
        x_i=\qty(b_i-\Sum[j=1][i-1]l_{ij}x_j)/l_{ii}
    \end{Equation}
\end{BoxFormula}

在数学之外，一个真正实用的算法需要考虑更多问题：如何利用矩阵的特点更高效的完成的形式上的矩阵乘法计算？如何尽可能的进行原位计算以节约空间消耗？例如，在\cref{fml:前向代换}的计算中，注意到$x_i$只依赖于$b_i$，因此，计算出的$\vb{x}$完全可以覆盖掉$\vb{b}$原本占用的存储空间。
\begin{BoxAlgorithm}[行主序的前向代换]
    若$\vb{L}\in\R^{n\times n}$是非奇异的下三角矩阵，且$\vb{b}\in\R^{n}$，该算法会以$\vb{L}\vb{x}=\vb{b}$的解覆写$\vb{b}$。
    \begin{Algoplain}
        $\vb{b}(1)=\vb{b}(1)/\vb{L}(1,1)$\;
        \For{$i=2:n$}{
            $\vb{b}(i)=[\vb{b}(i)-\vb{L}(i,1:i-1)\cdot\vb{b}(1:i-1)]/\vb{L}(i,i)$\;
        }
    \end{Algoplain}
\end{BoxAlgorithm}

\subsection{反向代换}

考虑一个$2\times 2$的上三角系统

\begin{Equation}[2x2反向代换1]
    \mqty[
        u_{11} & u_{12} \\
        0      & u_{22} \\
    ]
    \mqty[
        x_1 \\
        x_2 \\
    ]
    =
    \mqty[
        b_1 \\
        b_2 \\
    ]
\end{Equation}

若$u_{11},u_{22}\neq 0$，则未知数可以从下到上依次确定
\begin{Gather}[2x2反向代换]
    x_2=b_2/u_{22}\id{2}\\
    x_1=(b_1-u_{12}x_2)/u_{11}\id{3}
\end{Gather}

这一问题的一般形式可以概括为
\begin{BoxFormula}[反向代换]
    若$\vb{U}\in\R^{n\times n}$是上三角矩阵，而$\vb{b}\in\R^n$，那么
    \begin{Equation}[上三角系统]
        \vb{U}\vb{x}=\vb{b}
    \end{Equation}

    该方程可以用反向代换（Backward Substitution）的方式求出
    \begin{Equation}[反向代换]
        x_i=\qty(b_i-\Sum[j=i+1][n]u_{ij}x_j)/u_{ii}
    \end{Equation}
\end{BoxFormula}

在算法层面，前向代换和反向代换是类似的
\begin{BoxAlgorithm}[行主序的反向代换]
    若$\vb{U}\in\R^{n\times n}$是非奇异的上三角矩阵，且$\vb{b}\in\R^{n}$，该算法会以$\vb{U}\vb{x}=\vb{b}$的解覆写$\vb{b}$。
    \begin{Algoplain}
        $\vb{b}(n)=\vb{b}(n)/\vb{U}(n,n)$\;
        \For{$i=n-1:-1:1$}{
            $\vb{b}(i)=[\vb{b}(i)-\vb{U}(i,i+1:n)\cdot\vb{b}(i+1:n)]/\vb{U}(i,i)$\;
        }
    \end{Algoplain}
\end{BoxAlgorithm}

关于两种代换的名称，在此做一个总结
\begin{itemize}
    \item 前向代换适用于下三角系统$\vb{L}\hspace{1.5pt}\vb{x}=\vb{b}$，因为$\vb{x}$是从第$1$行到第$n$行定出的。
    \item 反向代换适用于上三角系统$\vb{U}\vb{x}=\vb{b}$，因为$\vb{x}$是从第$n$行到第$1$行定出的。
\end{itemize}

\subsection{列主序的三角系统解法}

在\cref{subsec:前向代换,subsec:反向代换}中，使用的都是行主序（Row Oriented）的求解算法，我们逐行求出各个未知数，并在计算每个新的未知数再移除所有已知的未知数。而本节会介绍另外一种成为列主序（Column Oriented）的求解算法，我们在计算出一个未知数后会立即将其所在列移除。

不妨考虑一个$3\times 3$的下三角系统$\vb{L}\vb{x}=\vb{b}$
\begin{Equation}[列主序三角系统1]
    \mqty[
        l_{11} & 0      & 0      \\
        l_{21} & l_{22} & 0      \\
        l_{31} & l_{32} & l_{33} \\
    ]
    \mqty[
        x_1 \\
        x_2 \\
        x_3 \\
    ]
    =
    \mqty[
        b_1 \\
        b_2 \\
        b_3 \\
    ]
\end{Equation}

显然$x_1=b_1/l_{11}$，随后就可以将三角系统的规模从$\R^{3\times 3}$降低到$\R^{2\times 2}$
\begin{Equation}[列主序三角系统2]
    \mqty[
        l_{22} & 0      \\
        l_{32} & l_{33} \\
    ]
    \mqty[
        x_2 \\
        x_3 \\
    ]
    =
    \mqty[
        b_2 \\
        b_3 \\
    ]
    -
    \mqty[
        l_{21} \\
        l_{31} \\
    ] 
    x_1
\end{Equation}

这就是列主序的想法，在计算量上和行主序没有区别，但给出了另一种可用的计算顺序。

\begin{BoxAlgorithm}[列主序的前向代换]
    若$\vb{L}\in\R^{n\times n}$是非奇异的下三角矩阵，且$\vb{b}\in\R^{n}$，该算法会以$\vb{L}\vb{x}=\vb{b}$的解覆写$\vb{b}$。
    \begin{Algoplain}
        \For{$j=1:n-1$}{
            $\vb{b}(j)=\vb{b}(j)/\vb{L}(j,j)$\;
            $\vb{b}(j+1:n)=\vb{b}(j+1:n)-\vb{b}(j)\cdot\vb{b}{L}(j+1:n,j)$\;
        }
        $\vb{b}(n)=\vb{b}/\vb{L}(n,n)$\;
    \end{Algoplain}
\end{BoxAlgorithm}

\begin{BoxAlgorithm}[列主序的反向代换]
    若$\vb{U}\in\R^{n\times n}$是非奇异的上三角矩阵，且$\vb{b}\in\R^{n}$，该算法会以$\vb{U}\vb{x}=\vb{b}$的解覆写$\vb{b}$。
    \begin{Algoplain}
        \For{$j=n:-1:2$}{
            $\vb{b}(j)=\vb{b}(j)/\vb{U}(j,j)$\;
            $\vb{b}(1:j-1)=\vb{b}(1:j-1)-\vb{b}(j)\cdot\vb{U}(1:j-1,j)$
        }
        $\vb{b}(1)=\vb{b}/\vb{L}(1,1)$\;
    \end{Algoplain}
\end{BoxAlgorithm}
